% Source: ME15_v1.html
\documentclass[11pt]{article}
\usepackage[margin=1in]{geometry}
\usepackage{amsmath}
\usepackage{amssymb}
\usepackage{siunitx}
\usepackage{enumitem}

\title{ME15: Waves}
\author{}
\date{}

\begin{document}
\maketitle

\section*{ME15: Waves}

\begin{enumerate}[leftmargin=*,label=\textbf{Question \arabic*.}]
\item \hfill \textit{/ 10 pts}

The wavelength of visible light ranges from about 380 − 750 nm. Calculate the frequency of visible light at 564 nm and fill in the missing factor below. The speed of light is \emph{c} = 2.9979 × 10$^{8}$ m/s.



\emph{f }= \rule{2cm}{0.4pt} × 10$^{12}$ Hz (Note the units!)

\textit{Correct answer: } 531.5426

\item \hfill \textit{/ 20 pts}

Suppose a wave has frequency 36.0 Hz and wavelength 4.20 m. Answer the following questions:

\textbf{(a) }The speed of the wave is [Select]
 m/s.

\textbf{(b)} The angular frequency is 226
 rad/s.

\textbf{(c)} The period is [Select]
 s.

\textbf{(d)} The wave number is [Select]
 rad/m.

\textit{Answer 1:}
\begin{itemize}
  \item 132
  \item 113
  \item 142
  \item 125
  \item \textbf{[correct]} 151
\end{itemize}
\textit{Answer 2:}
\begin{itemize}
  \item 208
  \item \textbf{[correct]} 226
  \item 113
  \item 238
  \item 300
\end{itemize}
\textit{Answer 3:}
\begin{itemize}
  \item \textbf{[correct]} 0.0278
  \item 0.0297
  \item 0.0254
  \item 0.0332
  \item 0.0413
\end{itemize}
\textit{Answer 4:}
\begin{itemize}
  \item 0.238
  \item 0.118
  \item 1.24
  \item 1.39
  \item 0.184
  \item 0.412
  \item 0.200
  \item 1.82
  \item \textbf{[correct]} 1.50
  \item 1.16
\end{itemize}

\item \hfill \textit{/ 10 pts}

Suppose a solution to the wave equation is given by \emph{y}(\emph{x},\emph{t}) = \emph{A} exp(𝝎\emph{t} + \emph{kx}), where \emph{A} = 5 m, 𝝎 = 20 rad/s and \emph{k} = 0.8 rad/m. If the wave equation is given by $\frac{\partial^2y}{\partial t^2}-\beta^2\frac{\partial^2y}{\partial x^2}=\:0$ , determine the magnitude of $\beta$ in MKS units.

\textit{Correct answer: } 25

\item \hfill \textit{/ 30 pts}

Suppose a wave function is given by \emph{y}(\emph{t},\emph{x}) = 5 cos(3\emph{t} + 0.5\emph{x} - 2.1) where all values are in MKS units. Then ...

\textbf{(a)} the amplitude is 5
 m;

\textbf{(b)}the angular frequency is [Select]
 rad/s;

\textbf{(c)} the wave number is [Select]
 rad/m;

\textbf{(d) }the phase shift is [Select]
 rad;

\textbf{(e) }the frequency is [Select]
 Hz;

\textbf{(f)} the period is 2.09
 s;

\textbf{(}\textbf{g)} the wavelength is [Select]
 m;

\textbf{(h)} the propagation speed of the wave is [Select]
 m/s;

\textbf{(i)} the transverse \textbf{speed} of the wave at (0,0) is [Select]
 m/s.

\textit{Answer 1:}
\begin{itemize}
  \item -2.1
  \item 3
  \item \textbf{[correct]} 5
  \item -5
  \item 0.5
\end{itemize}
\textit{Answer 2:}
\begin{itemize}
  \item 0.333
  \item -2.1
  \item 5
  \item \textbf{[correct]} 3
  \item 0.5
\end{itemize}
\textit{Answer 3:}
\begin{itemize}
  \item 2
  \item -2.1
  \item \textbf{[correct]} 0.5
  \item 1.5
  \item 3
\end{itemize}
\textit{Answer 4:}
\begin{itemize}
  \item 3
  \item -0.7
  \item \textbf{[correct]} -2.1
  \item 0.5
  \item 5
\end{itemize}
\textit{Answer 5:}
\begin{itemize}
  \item -2.1
  \item \textbf{[correct]} 0.477
  \item 0.121
  \item 3
  \item 2.09
  \item 0.5
  \item 0.0796
\end{itemize}
\textit{Answer 6:}
\begin{itemize}
  \item 0.729
  \item \textbf{[correct]} 2.09
  \item 3
  \item 0.477
  \item -2.1
\end{itemize}
\textit{Answer 7:}
\begin{itemize}
  \item 0.477
  \item 8.25
  \item \textbf{[correct]} 12.6
  \item -2.1
  \item 5
\end{itemize}
\textit{Answer 8:}
\begin{itemize}
  \item 5
  \item \textbf{[correct]} 6.0
  \item 6.4
  \item 4.6
  \item 5.7
\end{itemize}
\textit{Answer 9:}
\begin{itemize}
  \item \textbf{[correct]} 12.9
  \item 13.7
  \item 6.01
  \item -6.83
  \item -8.73
\end{itemize}

\item \hfill \textit{/ 10 pts}

A wire of mass 12 g and length 0.7 m under a tension of 514 N has an amplitude of 0.8 cm. If the frequency is 123 Hz, calculate the average power in watts of a wave propagating on the wire.



\emph{P}$_{av}$ = \rule{2cm}{0.4pt} W

\textit{Correct answer: } 56.734

\item \hfill \textit{/ 10 pts}

Two traveling waves of the same frequency (\emph{f} = 200 Hz) and wavelength interfere to produce a standing wave. The distance from node \#1 to node \#6 is 4.89 m. What is the wavelength (in m) of the original waves?

\textit{Correct answer: } 1.956

\item \hfill \textit{/ 10 pts}

A vibrating string has fourth harmonic frequency of 242 Hz. What is the third harmonic in Hertz?

\textit{Correct answer: } 182

\item \hfill \textit{/ 10 pts}

A wire with mass 0.066 kg and length 0.4 m has a fundamental frequency of 107 Hz. Calculate the tension in the wire in Newtons.

\textit{Correct answer: } 1,209

\item \hfill \textit{/ 10 pts}

A wire with mass 6 g is under a tension of 759 N. If the wire has fundamental frequency 318 Hz, find the length of the wire. (Be careful with units!)

\textit{Correct answer: } 0.3127

\item \hfill \textit{/ 10 pts}

Using only words from the list below, place the correct or best word in each blank in the following sentences concerning standing waves on a string fixed at both ends. Words may be used more than once.

The distance from a node to an adjacent antinode is
 of a wavelength. The distance from one end of the string to the other is
 of a wavelength at the fundamental frequency.

Word list:\textbf{ one-eighth; one-quarter; three-eighths; one-third; one-half; five-eighths; two-thirds; three-quarters}

\textit{Answer 1:}
\begin{itemize}
  \item \textbf{[correct]} one-quarter
\end{itemize}
\textit{Answer 2:}
\begin{itemize}
  \item \textbf{[correct]} one-half
\end{itemize}

\item \hfill \textit{/ 10 pts}

A wave whereby the vibration is perpendicular to the direction of propagation is called a transverse
 wave, whereas a wave having the vibration parallel to the direction of motion is called a longitudinal
 wave.

\textit{Answer 1:}
\begin{itemize}
  \item \textbf{[correct]} transverse
  \item longitudinal
\end{itemize}
\textit{Answer 2:}
\begin{itemize}
  \item \textbf{[correct]} longitudinal
  \item transverse
\end{itemize}

\item \hfill \textit{/ 10 pts}

Using only words from the list below, place the correct or best word in each blank in the following sentences concerning waves on a vibrating guitar string. Words may be used more than once.

Doubling the tension of a guitar string
 the speed of the wave on the string and
 the frequency of vibration. Using a string of four times the mass
 the speed of the wave by a factor of
 and
 the frequency by the same factor.



Word List: \textbf{doubles; increases; decreases; one-half; one-third; one-fourth; triples; quadruples}

\textit{Answer 1:}
\begin{itemize}
  \item \textbf{[correct]} increases
\end{itemize}
\textit{Answer 2:}
\begin{itemize}
  \item \textbf{[correct]} increases
\end{itemize}
\textit{Answer 3:}
\begin{itemize}
  \item \textbf{[correct]} decreases
\end{itemize}
\textit{Answer 4:}
\begin{itemize}
  \item \textbf{[correct]} one-half
\end{itemize}
\textit{Answer 5:}
\begin{itemize}
  \item \textbf{[correct]} decreases
\end{itemize}

\item \hfill \textit{/ 10 pts}

A given wave will have a certain number of wavelengths per meter. The number of radians corresponding to that number of wavelengths is called the [blank].

\begin{itemize}
  \item \textbf{[correct]} angular wavenumber
  \item \textbf{[correct]} wave number
  \item \textbf{[correct]} angular wave number
  \item \textbf{[correct]} wavenumber
\end{itemize}

\end{enumerate}

\end{document}
